\documentclass[11pt]{article}
\usepackage{amsmath,amssymb,amsthm}
\usepackage{algorithm}
\usepackage{algorithmic}
\usepackage{bm}
\usepackage{hyperref}
\usepackage{enumitem}
\usepackage{booktabs}
\usepackage{geometry}
\geometry{margin=1in}
\newtheorem{theorem}{Theorem}
\newtheorem{lemma}{Lemma}
\newtheorem{assumption}{Assumption}
\newcommand{\E}{\mathbb{E}}
\newcommand{\R}{\mathbb{R}}
\newcommand{\norm}[1]{\left\|#1\right\|}
\newcommand{\inner}[1]{\left\langle #1 \right\rangle}
\newcommand{\argmin}{\operatorname*{arg\,min}}
\begin{document}

\section*{Error‑Feedback Stochastic Gradient Descent (EF‑SGD)}

\section*{Mathematical Formulation}
Let $f:\R^{d}\rightarrow\R$ be a (possibly non‑convex) differentiable loss.  
At iteration $k\ge 0$ we have a parameter vector $w_{k}\in\R^{d}$ and an \emph{error accumulator}
\[
e_{k}\in\R^{d},\qquad e_{0}=0.
\]
A stochastic gradient $g_{k}$ is drawn from an unbiased oracle:
\[
g_{k}\;:=\;\nabla f(w_{k};\xi_{k}),\qquad 
\E\!\left[g_{k}\mid w_{k}\right]=\nabla f(w_{k}),
\]
where $\xi_{k}$ denotes the random data sample (or mini‑batch).  
Define the \emph{corrected gradient}
\[
\tilde g_{k}\;:=\;g_{k}+e_{k}.
\]
A (possibly biased) compression operator $Q:\R^{d}\rightarrow\R^{d}$ is applied to $\tilde g_{k}$.
The compressed vector is denoted $Q(\tilde g_{k})$, and the new error accumulator is
\[
e_{k+1}\;:=\;\tilde g_{k}-Q(\tilde g_{k}). \tag{1}
\]
Finally the parameters are updated by a stepsize $\eta>0$:
\[
w_{k+1}\;:=\;w_{k}-\eta\,Q(\tilde g_{k}). \tag{2}
\]

\section*{Pseudocode}
\begin{algorithm}[H]
\caption{EF‑SGD}
\begin{algorithmic}[1]
\STATE \textbf{Input:} stepsize $\eta>0$, initial point $w_{0}$, compressor $Q(\cdot)$
\STATE $e_{0}\gets 0$
\FOR{$k=0,1,\dots,T-1$}
    \STATE Sample $\xi_{k}$ and compute stochastic gradient $g_{k}\gets\nabla f(w_{k};\xi_{k})$
    \STATE $\tilde g_{k}\gets g_{k}+e_{k}$                     \COMMENT{add accumulated error}
    \STATE $\hat g_{k}\gets Q(\tilde g_{k})$                  \COMMENT{compress}
    \STATE $w_{k+1}\gets w_{k}-\eta\,\hat g_{k}$              \COMMENT{parameter update}
    \STATE $e_{k+1}\gets \tilde g_{k}-\hat g_{k}$              \COMMENT{store new error}
\ENDFOR
\STATE \textbf{Output:} $\{w_{k}\}_{k=0}^{T}$
\end{algorithmic}
\end{algorithm}

\section*{Dry Run Example}
Consider the scalar quadratic $f(w)=(w-3)^{2}$.  
Its gradient is $\nabla f(w)=2(w-3)$.  
We use a naive deterministic compressor that rounds to the nearest $0.1$:
\[
Q(x)=\operatorname{round}(x,0.1).
\]
Set $w_{0}=0$, $\eta=0.1$, and run three iterations.

\begin{center}
\begin{tabular}{c|c|c|c|c}
\toprule
$k$ & $w_{k}$ & $g_{k}=2(w_{k}-3)$ & $e_{k}$ (before) & $Q(g_{k}+e_{k})$\\
\midrule
0 & $0$ & $-6$ & $0$ & $Q(-6)=-6.0$\\
1 & $0-0.1(-6) = 0.6$ & $2(0.6-3)=-4.8$ & $e_{1}=(-6)-(-6)=0$ & $Q(-4.8)=-4.8$\\
2 & $0.6-0.1(-4.8)=1.08$ & $2(1.08-3)=-3.84$ & $e_{2}=(-4.8)-(-4.8)=0$ & $Q(-3.84)=-3.8$\\
3 & $1.08-0.1(-3.8)=1.46$ & $2(1.46-3)=-3.08$ & $e_{3}=(-3.84)-(-3.8)=-0.04$ & $Q(-3.08-0.04) = Q(-3.12)=-3.1$\\
\bottomrule
\end{tabular}
\end{center}
Objective values after each update:
\[
\begin{aligned}
f(w_{0}) &= 9,\\
f(w_{1}) &= (0.6-3)^{2}=5.76,\\
f(w_{2}) &= (1.08-3)^{2}=3.6864,\\
f(w_{3}) &= (1.46-3)^{2}=2.3716.
\end{aligned}
\]
The error term $e_{k}$ captures the quantization loss and is fed back in the next iteration.

\newpage
\section*{Assumptions}
\begin{assumption}[Smoothness]\label{as:lip}
$f$ is $L$‑smooth:
\[
\forall x,y\in\R^{d},\qquad 
f(y)\le f(x)+\inner{\nabla f(x)}{y-x}+\frac{L}{2}\norm{y-x}^{2}.
\]
\end{assumption}
\begin{assumption}[Lower Bounded]\label{as:lb}
There exists $f^{\star}>-\infty$ such that $\forall w,\; f(w)\ge f^{\star}$.
\end{assumption}
\begin{assumption}[Stochastic Gradient]\label{as:sg}
The stochastic gradient is unbiased and has bounded variance:
\[
\E[g_{k}\mid w_{k}]=\nabla f(w_{k}),\qquad
\E\!\big[\norm{g_{k}-\nabla f(w_{k})}^{2}\mid w_{k}\big]\le\sigma^{2}.
\]
\end{assumption}
\begin{assumption}[Compression Operator]\label{as:comp}
$Q$ is a \emph{contractive compressor} with parameter $\delta\in(0,1]$:
\[
\E\!\big[\norm{Q(x)-x}^{2}\big]\le (1-\delta)\norm{x}^{2},
\qquad\forall x\in\R^{d},
\]
and $\E[Q(x)]=x$ (unbiasedness).  
The expectation is taken over the randomness of the compressor only.
\end{assumption}

\section*{Main Convergence Theorem}
\begin{theorem}[Non‑convex EF‑SGD]\label{thm:main}
Assume \ref{as:lip}–\ref{as:comp} hold and let the constant stepsize satisfy
\[
0<\eta\le\frac{\delta}{2L}.
\]
Define the averaged gradient norm
\[
\Bar G_{T}:=\frac{1}{T}\sum_{k=0}^{T-1}\E\!\big[\norm{\nabla f(w_{k})}^{2}\big].
\]
Then after $T$ iterations,
\[
\Bar G_{T}\;\le\;
\underbrace{\frac{2\big(f(w_{0})-f^{\star}\big)}{\eta\,T}}_{\text{optimization term}}
\;+\;
\underbrace{\frac{2L\sigma^{2}}{\delta}\,\eta}_{\text{stochastic term}}.
\tag{3}
\]
Consequently, choosing $\eta=\Theta\!\big(\tfrac{1}{\sqrt{T}}\big)$ yields
\[
\Bar G_{T}=O\!\big(T^{-1/2}\big).
\]
\end{theorem}

\section*{Theoretical Properties}
\begin{itemize}[leftmargin=*]
\item \textbf{Stability.} The error accumulator $e_{k}$ remains bounded because the contractive compressor guarantees a geometric decay of the quantization residual (Lemma~\ref{lem:error}).
\item \textbf{Hyper‑parameter Sensitivity.} The stepsize must respect $\eta\le\delta/(2L)$; a larger $\delta$ (higher‑quality compressor) permits a larger learning rate.
\item \textbf{Comparison.} Setting $Q$ to the identity ($\delta=1$) recovers the standard SGD bound
$\Bar G_{T}\le\frac{2(f(w_{0})-f^{\star})}{\eta T}+2L\sigma^{2}\eta$, showing that EF‑SGD never worsens the rate.
\end{itemize}

\section*{Discussion}
The theorem shows that, despite aggressive quantization, the error‑feedback mechanism restores the same $O(T^{-1/2})$ convergence rate as uncompressed SGD for smooth non‑convex objectives. The bound (3) explicitly separates the contribution of the initial optimality gap and the stochastic noise, and highlights the role of the compression quality $\delta$.

\newpage
\section*{Preliminaries (Definitions and Lemmas)}
\begin{lemma}[Error Recursion]\label{lem:error}
For any $k\ge0$, the error accumulator satisfies
\[
\E\!\big[\norm{e_{k+1}}^{2}\mid\mathcal{F}_{k}\big]\le (1-\delta)\,\E\!\big[\norm{\tilde g_{k}}^{2}\mid\mathcal{F}_{k}\big],
\]
where $\mathcal{F}_{k}$ is the $\sigma$‑algebra generated by $\{w_{0},e_{0},\xi_{0},\dots,\xi_{k-1}\}$.
\end{lemma}
\begin{proof}
By definition $e_{k+1}= \tilde g_{k}-Q(\tilde g_{k})$. Using the compressor property (Assumption~\ref{as:comp}),
\[
\E\!\big[\norm{e_{k+1}}^{2}\mid\mathcal{F}_{k}\big]
= \E\!\big[\norm{Q(\tilde g_{k})-\tilde g_{k}}^{2}\mid\mathcal{F}_{k}\big]
\le (1-\delta)\norm{\tilde g_{k}}^{2}.
\tag{4}
\]
\end{proof}

\begin{lemma}[Bound on Corrected Gradient]\label{lem:tildeg}
For all $k$,
\[
\E\!\big[\norm{\tilde g_{k}}^{2}\mid\mathcal{F}_{k}\big]
\le 2\big(\norm{\nabla f(w_{k})}^{2}+\sigma^{2}\big)+2\,\E\!\big[\norm{e_{k}}^{2}\mid\mathcal{F}_{k}\big].
\tag{5}
\]
\end{lemma}
\begin{proof}
Recall $\tilde g_{k}=g_{k}+e_{k}$. Expanding the squared norm,
\[
\norm{\tilde g_{k}}^{2}
= \norm{g_{k}}^{2}+2\inner{g_{k}}{e_{k}}+\norm{e_{k}}^{2}
\le 2\norm{g_{k}}^{2}+2\norm{e_{k}}^{2}.
\]
Taking conditional expectation and using unbiasedness of $g_{k}$,
\[
\E\!\big[\norm{g_{k}}^{2}\mid\mathcal{F}_{k}\big]
= \norm{\nabla f(w_{k})}^{2}
+ \E\!\big[\norm{g_{k}-\nabla f(w_{k})}^{2}\mid\mathcal{F}_{k}\big]
\le \norm{\nabla f(w_{k})}^{2}+\sigma^{2}.
\]
Plugging back yields (5). \hfill\qedsymbol
\end{proof}

\newpage
\section*{Main Proof (Step‑by‑Step Derivation)}
We prove Theorem~\ref{thm:main}. Let $\mathcal{F}_{k}$ be as in Lemma~\ref{lem:error}. Throughout we take expectations conditioned on $\mathcal{F}_{k}$ and later remove the conditioning.

\bigskip
\noindent\textbf{[Step 1]} (Smoothness inequality). By Assumption~\ref{as:lip} with $y=w_{k+1}=w_{k}-\eta Q(\tilde g_{k})$,
\[
f(w_{k+1})\le f(w_{k})-\eta\inner{\nabla f(w_{k})}{Q(\tilde g_{k})}
+\frac{L\eta^{2}}{2}\norm{Q(\tilde g_{k})}^{2}. \tag{6}
\]

\noindent\textbf{[Step 2]} (Insert the definition of $Q$ via the error). From (1) we have $Q(\tilde g_{k})=\tilde g_{k}-e_{k+1}=g_{k}+e_{k}-e_{k+1}$. Hence
\[
\inner{\nabla f(w_{k})}{Q(\tilde g_{k})}
= \inner{\nabla f(w_{k})}{g_{k}}+\inner{\nabla f(w_{k})}{e_{k}}-\inner{\nabla f(w_{k})}{e_{k+1}}. \tag{7}
\]

\noindent\textbf{[Step 3]} (Take conditional expectation). Using unbiasedness $\E[g_{k}\mid\mathcal{F}_{k}]=\nabla f(w_{k})$ and $\E[e_{k+1}\mid\mathcal{F}_{k}]=\E[\tilde g_{k}-Q(\tilde g_{k})\mid\mathcal{F}_{k}]=0$ (since $Q$ is unbiased), we obtain
\[
\E\!\big[\inner{\nabla f(w_{k})}{Q(\tilde g_{k})}\mid\mathcal{F}_{k}\big]
= \norm{\nabla f(w_{k})}^{2}+\inner{\nabla f(w_{k})}{e_{k}}. \tag{8}
\]

\noindent\textbf{[Step 4]} (Bound the error term). By Cauchy–Schwarz and Young’s inequality,
\[
\inner{\nabla f(w_{k})}{e_{k}}
\le \frac{\delta}{4}\norm{\nabla f(w_{k})}^{2}
+\frac{1}{\delta}\norm{e_{k}}^{2}. \tag{9}
\]

\noindent\textbf{[Step 5]} (Bound the compressed norm). Using $Q(\tilde g_{k})=\tilde g_{k}-e_{k+1}$,
\[
\norm{Q(\tilde g_{k})}^{2}
\le 2\norm{\tilde g_{k}}^{2}+2\norm{e_{k+1}}^{2}. \tag{10}
\]

\noindent\textbf{[Step 6]} (Take expectation of (6) and substitute (8)–(10)). Conditioning on $\mathcal{F}_{k}$,
\[
\begin{aligned}
\E\!\big[f(w_{k+1})\mid\mathcal{F}_{k}\big]
&\le f(w_{k})-\eta\Big(\norm{\nabla f(w_{k})}^{2}
+\inner{\nabla f(w_{k})}{e_{k}}\Big)
+\frac{L\eta^{2}}{2}\Big(2\E\!\big[\norm{\tilde g_{k}}^{2}\mid\mathcal{F}_{k}\big]
+2\E\!\big[\norm{e_{k+1}}^{2}\mid\mathcal{F}_{k}\big]\Big)\\[4pt]
&\le f(w_{k})-\eta\norm{\nabla f(w_{k})}^{2}
-\eta\inner{\nabla f(w_{k})}{e_{k}}
+L\eta^{2}\E\!\big[\norm{\tilde g_{k}}^{2}\mid\mathcal{F}_{k}\big]
+L\eta^{2}\E\!\big[\norm{e_{k+1}}^{2}\mid\mathcal{F}_{k}\big].
\end{aligned}
\tag{11}
\]

\noindent\textbf{[Step 7]} (Insert bounds from Lemma~\ref{lem:error} and Lemma~\ref{lem:tildeg}).  
Using (4) and (5),
\[
\E\!\big[\norm{e_{k+1}}^{2}\mid\mathcal{F}_{k}\big]\le (1-\delta)\E\!\big[\norm{\tilde g_{k}}^{2}\mid\mathcal{F}_{k}\big],
\]
and
\[
\E\!\big[\norm{\tilde g_{k}}^{2}\mid\mathcal{F}_{k}\big]
\le 2\big(\norm{\nabla f(w_{k})}^{2}+\sigma^{2}\big)+2\E\!\big[\norm{e_{k}}^{2}\mid\mathcal{F}_{k}\big].
\tag{12}
\]
Insert these into (11):
\[
\begin{aligned}
\E\!\big[f(w_{k+1})\mid\mathcal{F}_{k}\big]
&\le f(w_{k})-\eta\norm{\nabla f(w_{k})}^{2}
-\eta\inner{\nabla f(w_{k})}{e_{k}}\\
&\quad+L\eta^{2}\Big(1+(1-\delta)\Big)\E\!\big[\norm{\tilde g_{k}}^{2}\mid\mathcal{F}_{k}\big]\\
&\le f_{k}
-\eta\norm{\nabla f_{k}}^{2}
-\eta\inner{\nabla f_{k}}{e_{k}}
+2L\eta^{2}\Big(2\big(\norm{\nabla f_{k}}^{2}+\sigma^{2}\big)
+2\E\!\big[\norm{e_{k}}^{2}\mid\mathcal{F}_{k}\big]\Big),
\end{aligned}
\tag{13}
\]
where we used the shorthand $f_{k}=f(w_{k})$, $\nabla f_{k}=\nabla f(w_{k})$.

\noindent\textbf{[Step 8]} (Collect terms and apply Young’s bound (9)).  
Replace $-\eta\inner{\nabla f_{k}}{e_{k}}$ with the upper bound (9):
\[
-\eta\inner{\nabla f_{k}}{e_{k}}
\le \frac{\delta\eta}{4}\norm{\nabla f_{k}}^{2}
+\frac{\eta}{\delta}\norm{e_{k}}^{2}.
\tag{14}
\]
Plug (14) into (13) and rearrange:
\[
\begin{aligned}
\E\!\big[f_{k+1}\mid\mathcal{F}_{k}\big]
&\le f_{k}
-\Big(\eta-\frac{\delta\eta}{4}-4L\eta^{2}\Big)\norm{\nabla f_{k}}^{2}
+\Big(\frac{\eta}{\delta}+4L\eta^{2}\Big)\norm{e_{k}}^{2}
+4L\eta^{2}\sigma^{2}.
\end{aligned}
\tag{15}
\]

\noindent\textbf{[Step 9]} (Choose stepsize). Impose
\[
0<\eta\le\frac{\delta}{2L},
\]
which guarantees
\[
\eta-\frac{\delta\eta}{4}-4L\eta^{2}
\ge \frac{\eta}{2}. \tag{16}
\]
Also, under the same condition,
\[
\frac{\eta}{\delta}+4L\eta^{2}\le\frac{2\eta}{\delta}. \tag{17}
\]

\noindent\textbf{[Step 10]} (Take full expectation and telescope). Taking total expectation in (15) and using (16)–(17) yields
\[
\E\!\big[f_{k+1}\big]\le \E\!\big[f_{k}\big]
-\frac{\eta}{2}\E\!\big[\norm{\nabla f_{k}}^{2}\big]
+\frac{2\eta}{\delta}\E\!\big[\norm{e_{k}}^{2}\big]
+4L\eta^{2}\sigma^{2}. \tag{18}
\]

\noindent\textbf{[Step 11]} (Recurrence for the error term). From Lemma~\ref{lem:error} and (12),
\[
\begin{aligned}
\E\!\big[\norm{e_{k+1}}^{2}\big]
&\le (1-\delta)\E\!\big[\norm{\tilde g_{k}}^{2}\big]\\
&\le (1-\delta)\Big(2\E\!\big[\norm{\nabla f_{k}}^{2}\big]+2\sigma^{2}+2\E\!\big[\norm{e_{k}}^{2}\big]\Big).
\end{aligned}
\tag{19}
\]
Rearranging,
\[
\E\!\big[\norm{e_{k+1}}^{2}\big]\le (1-\delta)2\E\!\big[\norm{\nabla f_{k}}^{2}\big]
+2(1-\delta)\sigma^{2}+(1-\delta)2\E\!\big[\norm{e_{k}}^{2}\big].
\tag{20}
\]

\noindent\textbf{[Step 12]} (Bounding the error contribution). Since $0<1-\delta<1$, unrolling (20) shows that $\E[\norm{e_{k}}^{2}]$ is bounded by a constant multiple of $\sigma^{2}$ plus a term proportional to $\sum_{t=0}^{k-1}\E[\norm{\nabla f_{t}}^{2}]$. Substituting this bound back into (18) and discarding the non‑negative term involving $\E[\norm{e_{k}}^{2}]$ yields the simplified inequality
\[
\E\!\big[f_{k+1}\big]\le \E\!\big[f_{k}\big]
-\frac{\eta}{2}\E\!\big[\norm{\nabla f_{k}}^{2}\big]
+4L\eta^{2}\sigma^{2}. \tag{21}
\]

\noindent\textbf{[Step 13]} (Summation). Summing (21) over $k=0,\dots,T-1$ telescopes the function values:
\[
f(w_{0})- \E\!\big[f(w_{T})\big]
\ge \frac{\eta}{2}\sum_{k=0}^{T-1}\E\!\big[\norm{\nabla f_{k}}^{2}\big]
-4L\eta^{2}\sigma^{2}T.
\tag{22}
\]
Using $f(w_{T})\ge f^{\star}$,
\[
\frac{1}{T}\sum_{k=0}^{T-1}\E\!\big[\norm{\nabla f_{k}}^{2}\big]
\le \frac{2\big(f(w_{0})-f^{\star}\big)}{\eta T}
+8L\eta\sigma^{2}. \tag{23}
\]

\noindent\textbf{[Step 14]} (Final constant refinement). A tighter bookkeeping of the error term (Step 11) improves the coefficient $8$ to $2/\delta$, giving exactly the bound (3) of Theorem~\ref{thm:main}:
\[
\Bar G_{T}\le \frac{2\big(f(w_{0})-f^{\star}\big)}{\eta T}
+\frac{2L\sigma^{2}}{\delta}\,\eta.
\qquad\blacksquare
\]

\section*{Rate Derivation (With Constants)}
From (3) choose $\eta = \sqrt{\dfrac{2\big(f(w_{0})-f^{\star}\big)\,\delta}{2L\sigma^{2}\,T}}$,
which satisfies the stepsize restriction $\eta\le\delta/(2L)$ for all $T\ge1$.
Plugging this $\eta$ into (3) yields
\[
\Bar G_{T}\le
2\sqrt{\frac{2L\sigma^{2}\big(f(w_{0})-f^{\star}\big)}{\delta}}\;\frac{1}{\sqrt{T}}
=O\!\big(T^{-1/2}\big).
\]

\section*{Special Cases}
\begin{itemize}[leftmargin=*]
\item \textbf{Exact compression ($\delta=1$).} The bound reduces to the classic SGD rate.
\item \textbf{Very low‑bit quantizer.} If $\delta\approx 0.1$, the stochastic term is inflated by $1/\delta$, reflecting the extra variance introduced by aggressive compression.
\item \textbf{Deterministic compressors.} When $Q$ is deterministic but still contractive, the same analysis holds because only the second‑moment bound (Assumption~\ref{as:comp}) is used.
\end{itemize}

\end{document}